\chapter{Architettura}\label{ch:architettura}

Questo capitolo descrive che cosa emula vnic e con quale modello: la corrispondenza tra gli elementi di una SmartNIC e quelli del tool (sezione~\ref{sec:modello}), la topologia di rete che vnic crea (sezione~\ref{sec:topologia}), il comportamento atteso della VM per ciascuna azione XDP e i modi per osservarlo (sezione~\ref{sec:azioni-xdp}), e infine le operazioni di una SmartNIC, come si esprimono in eBPF e in P4 e quali sono implementate in vnic (sezione~\ref{sec:punti-comuni}).

\section{Modello: NIC con offload XDP}\label{sec:modello}

Tra le varie funzionalità delle SmartNIC, vnic emula in particolare l'offload di programmi XDP, cioè la possibilità di eseguire sulla scheda un programma eBPF che elabora i pacchetti in ricezione. Nel resto della tesi, con SmartNIC si intende quindi una scheda con questa funzionalità.

La SmartNIC vera e propria è rappresentata da una macchina virtuale Debian (indicata d'ora in poi come VM). La porta di rete della scheda corrisponde all'interfaccia \texttt{ens3} della VM, mentre il collegamento con l'host, che in una scheda reale è il bus PCIe verso il server, è rappresentato dall'interfaccia \texttt{ens4}. Il percorso interno alla scheda, dalla porta di rete verso l'host, è rappresentato dal bridge software \texttt{brVM}, che collega le due interfacce. Il programma offloaded, che sulla scheda reale elabora i pacchetti appena arrivano dalla rete, viene caricato su \texttt{ens3}: per ogni pacchetto decide cosa farne e, se lo lascia passare, il pacchetto raggiunge l'host attraverso \texttt{ens4}. Infine la rete esterna e il server che monta la scheda sono simulati da due network namespace: \texttt{ns0}, che invia pacchetti verso \texttt{ens3}, e \texttt{ns1}, che riceve da \texttt{ens4} i pacchetti filtrati.

Questo modello è volutamente semplice, in quanto deve servire a testare programmi XDP come su una vera SmartNIC. Dal suo utilizzo non è quindi possibile trarre conclusioni sulle prestazioni di una vera SmartNIC né sulle sue funzionalità di offload hardware, come il calcolo dei checksum o la segmentazione dei pacchetti TCP, che nell'emulazione vengono anzi disattivate quando si carica un programma XDP (capitolo~\ref{ch:valutazione}). La tabella~\ref{tab:modello} riassume la corrispondenza.

\begin{table}[H]
    \centering
    \caption{Corrispondenza tra gli elementi di una SmartNIC e quelli di vnic}
    \label{tab:modello}
    \begin{tabular}{@{}ll@{}}
        \toprule
        \textbf{SmartNIC} & \textbf{vnic} \\
        \midrule
        scheda di rete & VM QEMU/KVM \\
        porta di rete & \texttt{ens3} \\
        programma XDP eseguito sulla scheda & programma caricato su \texttt{ens3} \\
        percorso interno e collegamento con l'host & \texttt{brVM} e \texttt{ens4} \\
        rete esterna & namespace \texttt{ns0} \\
        server che monta la scheda & namespace \texttt{ns1} \\
        \bottomrule
    \end{tabular}
\end{table}

\section{Topologia}\label{sec:topologia}

La topologia di rete è visibile nella figura~\ref{fig:architettura}.

\begin{figure}[htb]
    \centering
    \begin{tikzpicture}[
    font=\small,
    iface/.style={draw, rounded corners=2pt, fill=white, minimum height=7mm, minimum width=15mm},
    xdp/.style={iface, fill=orange!25, draw=orange!70!black, very thick},
    zone/.style={draw, rounded corners=4pt},
    lbl/.style={font=\footnotesize\bfseries, anchor=north west, inner sep=3pt},
    note/.style={font=\scriptsize, text=black!70},
    tubo/.style={draw=black!60, line width=0.7pt, double=gray!30, double distance=4pt},
    bridge/.style={draw=blue!55!black, line width=1.5pt},
    nomebr/.style={font=\footnotesize\bfseries, text=blue!55!black, above, inner sep=2pt},
    qemu/.style={draw=black!70, thick},
]
    \node[xdp]   (ens3) at (3.2, 3.4) {ens3};
    \node[iface] (ens4) at (8.8, 3.4) {ens4};
    \node[note, above=0.5mm of ens3, text=orange!50!black] {programma XDP};
    \draw[bridge] (ens3) -- node[nomebr] {brVM} (ens4);
    \begin{scope}[on background layer]
        \node[zone, fill=gray!12, fit=(ens3)(ens4), inner xsep=8mm, inner ysep=7mm] (vm) {};
    \end{scope}
    \node[lbl, anchor=south west] at (vm.north west) {VM QEMU/KVM};

    \node[iface] (vb0)  at (0.2, 1.0) {veth0-br};
    \node[iface] (tap0) at (3.2, 1.0) {tap0};
    \node[iface] (tap1) at (8.8, 1.0) {tap1};
    \node[iface] (vb1)  at (11.8,1.0) {veth1-br};
    \draw[bridge] (vb0) -- node[nomebr] {br0} (tap0);
    \draw[bridge] (tap1) -- node[nomebr] {br1} (vb1);
    \begin{scope}[on background layer]
        \node[zone, dashed, fill=blue!4, fit=(vb0)(tap0)(tap1)(vb1), inner sep=3mm] (host) {};
    \end{scope}
    \node[lbl, anchor=south] at (host.south) {host};

    \node[iface] (v0) at (0.2,-1.9) {veth0};
    \node[iface] (v1) at (11.8,-1.9) {veth1};
    \begin{scope}[on background layer]
        \node[zone, fill=green!10, fit=(v0), inner xsep=6mm, inner ysep=5mm] (ns0) {};
        \node[zone, fill=green!10, fit=(v1), inner xsep=6mm, inner ysep=5mm] (ns1) {};
    \end{scope}
    \node[lbl] at (ns0.north west) {ns0};
    \node[lbl] at (ns1.north west) {ns1};
    \node[note, below=0.5mm of ns0] {192.168.100.1};
    \node[note, below=0.5mm of ns1] {192.168.100.2};

    \draw[tubo] (v0) -- node[note, right=2mm] {veth pair} (vb0);
    \draw[tubo] (v1) -- node[note, left=2mm] {veth pair} (vb1);
    \draw[qemu] (tap0) -- node[note, left] {virtio-net} (ens3);
    \draw[qemu] (tap1) -- node[note, right] {virtio-net} (ens4);
\end{tikzpicture}

    \caption{Topologia creata da vnic: la VM emula una SmartNIC con un programma XDP in ingresso}
    \label{fig:architettura}
\end{figure}

Un pacchetto viene inviato dal namespace \texttt{ns0} attraverso l'interfaccia \texttt{veth0}, con indirizzo IP 192.168.100.1/24. Questa fa parte di una veth pair, cioè una coppia di interfacce virtuali collegate tra loro come da un cavo, il cui altro capo, \texttt{veth0-br}, si trova nel namespace di default, quello dell'host. Qui il bridge software \texttt{br0} collega \texttt{veth0-br} all'interfaccia TAP \texttt{tap0}. Una TAP è un'interfaccia virtuale il cui altro capo non è un cavo ma un programma in spazio utente: i frame che l'host trasmette su \texttt{tap0} vengono letti da QEMU, che li consegna alla scheda di rete virtio-net emulata, vista dalla VM come \texttt{ens3}. Su \texttt{ens3} è stato caricato al boot il programma XDP, che elabora il pacchetto prima che entri nello stack di rete della VM. Se il programma lo lascia passare, il pacchetto raggiunge il bridge \texttt{brVM}, anch'esso creato al boot, che lo inoltra all'altra interfaccia della VM, \texttt{ens4}. Le altre azioni possibili sono descritte nella sezione~\ref{sec:azioni-xdp}.

Da qui il percorso è speculare a quello visto finora: QEMU scrive su \texttt{tap1} i frame trasmessi da \texttt{ens4}, e il bridge \texttt{br1} e la veth pair formata da \texttt{veth1-br} e \texttt{veth1}, con indirizzo IP 192.168.100.2/24, li portano nel namespace \texttt{ns1}.

I tre bridge (\texttt{br0}, \texttt{brVM} e \texttt{br1}) lavorano a livello 2 dello stack di rete e inoltrano i frame in base al MAC di destinazione, che è quello di \texttt{veth1}. Per \texttt{ns0} e \texttt{ns1} è quindi come essere collegati direttamente da un cavo, e infatti i loro indirizzi appartengono alla stessa sottorete. Anche i MAC delle interfacce della VM, che vnic fissa a \texttt{52:54:00:12:34:56} per \texttt{ens3} e a \texttt{52:54:00:12:34:57} per \texttt{ens4} (con il prefisso usato per convenzione da QEMU), non intervengono nell'inoltro.

Poiché XDP agisce solo sui pacchetti in ingresso, il programma viene eseguito soltanto nel verso da \texttt{ns0} a \texttt{ns1}. Il traffico di ritorno, per esempio le risposte ai ping, entra nella VM da \texttt{ens4}, dove non è caricato alcun programma, e raggiunge \texttt{ns0} senza essere filtrato.

\section{Azioni XDP nella VM}\label{sec:azioni-xdp}

A seconda dell'azione XDP restituita dal programma caricato, il risultato può variare. Per verificarlo è stato utilizzato un semplice programma eBPF di prova, \texttt{xdp\_actions.c} (listing~\ref{lst:xdp-actions}), che sceglie l'azione in base alla porta UDP di destinazione. Da \texttt{ns0} sono stati inviati 5 pacchetti per ciascuna porta, osservando cosa accade dentro e fuori dalla VM con i metodi descritti di seguito.

Per contare i pacchetti per ogni azione esistono 3 modi:
\begin{enumerate}
    \item direttamente nel programma XDP, con una mappa \texttt{BPF\_MAP\_TYPE\_PERCPU\_ARRAY} indicizzata per azione~\cite{xdp-tutorial}. Questo metodo è il più completo, perché è il programma stesso a sapere quale azione restituisce: distingue tutte e 5 le azioni, e volendo può contarle anche per porta o per flusso. Richiede però uno strumento per leggere la mappa, come \texttt{bpftool}, che nella VM di vnic non è presente;
    \item nel driver della VM, con i contatori mostrati da \texttt{ethtool -S ens3} e con i tracepoint \texttt{xdp:*} del kernel. TX e REDIRECT hanno un contatore dedicato, mentre DROP e ABORTED finiscono nello stesso contatore (\texttt{xdp\_drops}) e li distingue solo il tracepoint \texttt{xdp\_exception}, emesso per ABORTED. PASS non ha un contatore suo e si ricava per differenza dal totale dei pacchetti elaborati;
    \item esternamente alla VM, con \texttt{tcpdump} sulle interfacce \texttt{tap0} e \texttt{tap1}. Da qui ogni azione lascia una traccia diversa: con PASS il pacchetto esce da \texttt{tap1} verso \texttt{ns1}; con TX torna indietro su \texttt{tap0} verso \texttt{ns0}; con REDIRECT verso \texttt{ens4} esce anch'esso da \texttt{tap1}, perché salta \texttt{brVM} ma arriva comunque a \texttt{ens4}; con DROP e ABORTED entra da \texttt{tap0} e non esce più da nessuna parte. Dall'esterno, quindi, DROP e ABORTED sono indistinguibili, e lo stesso vale per PASS e REDIRECT riuscito.
\end{enumerate}

In vnic è possibile utilizzare direttamente il terzo metodo, ed è proprio quello che fa lo script \texttt{verify.sh} nei due esempi. Va notato che questo è anche l'unico disponibile per una SmartNIC vera vista dall'host senza gli strumenti del produttore: è quindi coerente con il modello di SmartNIC adottato.

Per l'esperimento, oltre al terzo metodo, è stato usato anche il secondo, così da confrontare ciò che si vede dall'esterno con ciò che registra il driver. A questo scopo è stata usata una copia di vnic modificata, che aggiunge le seguenti funzionalità:
\begin{itemize}
    \item attivazione dei tracepoint \texttt{xdp:*}, che il kernel emette per esempio quando un programma restituisce ABORTED o quando una redirect fallisce;
    \item stampa ogni secondo dell'output di \texttt{ethtool -S ens3} ed \texttt{ethtool -S ens4}, insieme ai tracepoint, sulla console seriale della VM, che QEMU salva in un file di log sull'host.
\end{itemize}

Nello stesso tempo \texttt{tcpdump} registra su \texttt{tap0} e \texttt{tap1} i frame in entrambe le direzioni. Per ogni porta si è poi confrontato il valore dei contatori prima e dopo l'invio dei 5 pacchetti.

I contatori mostrati da \texttt{ethtool -S} sono statistiche che il driver \texttt{virtio\_net} tiene per ogni coda dell'interfaccia. Poiché le interfacce della VM hanno una sola coda, i nomi iniziano con \texttt{rx\_queue\_0\_} per la ricezione e con \texttt{tx\_queue\_0\_} per la trasmissione. Quelli relativi a XDP sono:
\begin{itemize}
    \item \texttt{xdp\_packets}: pacchetti elaborati dal programma, qualunque sia l'azione;
    \item \texttt{xdp\_drops}: pacchetti scartati, sia con DROP che con ABORTED;
    \item \texttt{xdp\_tx}: pacchetti rimandati indietro con TX, contati sia in ricezione che in trasmissione;
    \item \texttt{xdp\_redirects}: pacchetti rediretti verso un'altra interfaccia.
\end{itemize}

Nella tabella~\ref{tab:azioni-xdp} sono riassunti i risultati delle azioni XDP in vnic.

% Risultati dell'esperimento sulle azioni XDP (sezione 3.3): 5 pacchetti per azione.
\begin{table}[H]
    \centering
    \footnotesize
    \caption{Azioni XDP nella VM: frame osservati su \texttt{tap0} e \texttt{tap1} e contatori di \texttt{ens3} che aumentano (5 pacchetti per azione)}
    \label{tab:azioni-xdp}
    \begin{tabular}{@{}lcc>{\raggedright\arraybackslash}p{5.6cm}@{}}
        \toprule
        \textbf{Azione} & \textbf{verso ns1} & \textbf{verso ns0} & \textbf{Contatori di \texttt{ens3}} \\
         & (\texttt{tap1}) & (\texttt{tap0}) & \\
        \midrule
        PASS                       & 5 & 0 & nessuno dedicato \\
        DROP                       & 0 & 0 & \texttt{xdp\_drops} \\
        ABORTED                    & 0 & 0 & \texttt{xdp\_drops}, tracepoint \texttt{xdp\_exception} \\
        TX                         & 0 & 5 & \texttt{xdp\_tx} \\
        REDIRECT, \texttt{ens4} senza XDP & 0 & 0 & \texttt{xdp\_redirects} \\
        REDIRECT, \texttt{ens4} con XDP   & 5 & 0 & \texttt{xdp\_redirects} \\
        \bottomrule
    \end{tabular}
\end{table}


Due risultati meritano attenzione. Il primo riguarda PASS: pur essendo l'azione più comune, è l'unica senza un contatore nel driver. Il secondo riguarda REDIRECT verso \texttt{ens4}, che funziona solo se anche \texttt{ens4} ha un programma XDP caricato. In caso contrario il driver conta comunque il pacchetto come rediretto, ma poi lo scarta al momento dell'invio. Dai contatori, quindi, un REDIRECT fallito sembra riuscito.

\section{Operazioni di una SmartNIC in eBPF, P4 e vnic}\label{sec:punti-comuni}

Il modello di elaborazione di una SmartNIC è lo stesso, indipendentemente dal linguaggio utilizzato. Per ogni pacchetto se ne analizzano gli header, si prende una decisione, eventualmente consultando tabelle o uno stato, e infine si applica un'azione, come scartarlo, consegnarlo all'host o inoltrarlo altrove. P4 è nato proprio per programmare questo tipo di elaborazione su hardware di rete~\cite{bosshart2014p4}, mentre un programma XDP può essere eseguito sia dal driver sia, con l'offload, direttamente dalla scheda~\cite{kicinski2016offload}. Il compilatore \texttt{p4c-ebpf} collega i due mondi, traducendo un programma P4 in un programma eBPF: in vnic, quindi, entrambi i linguaggi diventano un programma XDP caricato su \texttt{ens3}.

La tabella~\ref{tab:punti-comuni} elenca le principali operazioni che una SmartNIC può svolgere su un pacchetto e, per ciascuna, se e come la si può esprimere in eBPF e in P4, usando per quest'ultimo il modello \texttt{ebpfFilter} di \texttt{p4c-ebpf}.

% Operazioni di una SmartNIC con offload XDP e come si esprimono in eBPF e in P4 (sezione 3.4).
\begin{table}[H]
    \centering
    \footnotesize
    \caption{Operazioni di una SmartNIC e come si esprimono in eBPF e in P4}
    \label{tab:punti-comuni}
    \begin{tabular}{@{}>{\raggedright\arraybackslash}p{3.7cm}>{\raggedright\arraybackslash}p{5.1cm}>{\raggedright\arraybackslash}p{5.1cm}@{}}
        \toprule
        \textbf{SmartNIC} & \textbf{eBPF} & \textbf{P4} \\
        \midrule
        ricezione da più porte e code (RSS) & presente: \texttt{ingress\_ifindex} e \texttt{rx\_queue\_index} indicano porta e coda di arrivo & assente \\
        \addlinespace
        analisi degli header (parsing) & presente: controllo in C sulla lunghezza prima di ogni lettura (richiesto dal verifier) & presente: blocco dichiarativo \texttt{parser} con \texttt{extract} e \texttt{select} \\
        \addlinespace
        classificazione dei pacchetti con tabelle e stato & presente: mappe BPF, accessibili anche dallo spazio utente & presente: tabelle match-action tradotte in mappe BPF \\
        \addlinespace
        drop del pacchetto o consegna all'host & presente: \texttt{XDP\_DROP} e \texttt{XDP\_PASS} & presente: parametro booleano di uscita false o true \\
        \addlinespace
        rinvio sulla porta di arrivo o redirezione su altra porta del pacchetto & presente: \texttt{XDP\_TX} e \texttt{XDP\_REDIRECT} & assente \\
        \addlinespace
        modifica del pacchetto & presente & assente \\
        \bottomrule
    \end{tabular}
\end{table}


Ne risulta che eBPF può esprimere tutte le operazioni della tabella, mentre P4 con il modello \texttt{ebpfFilter} si limita a un filtro: analizza il pacchetto e decide se accettarlo o scartarlo, senza poterlo modificare, rimandare indietro o inoltrare altrove. In cambio, il parsing e le tabelle si descrivono in modo dichiarativo, senza dover gestire a mano i controlli sulla lunghezza del pacchetto.

Di queste operazioni vnic modella l'elaborazione del programma sulla porta di ingresso. Con entrambi i linguaggi il parsing e la decisione avvengono come su una SmartNIC reale. Con eBPF sono disponibili tutte le azioni, con il limite già visto per REDIRECT verso \texttt{ens4} (sezione~\ref{sec:azioni-xdp}); con P4, invece, solo quelle consentite dal modello \texttt{ebpfFilter}.

Alcuni aspetti, invece, non sono modellati. Il più significativo riguarda le code di ricezione. Una scheda reale ha più code e distribuisce i pacchetti tra i core con l'RSS, e il programma XDP può sapere da quale coda arriva ogni pacchetto. In vnic ogni interfaccia virtio ha una sola coda (lo mostra \texttt{ethtool -l ens3}), quindi \texttt{rx\_queue\_index} vale sempre 0 e questa dimensione del modello non è rappresentata.

Manca inoltre un piano di controllo. Le mappe e le tabelle vengono create quando il programma è caricato, ma dall'host non è possibile leggerle né modificarle: non si possono, per esempio, consultare dei contatori o aggiungere regole a una tabella mentre il programma è in esecuzione.

Infine vnic non modella gli offload hardware, come il calcolo dei checksum o la segmentazione dei pacchetti TCP, che su una SmartNIC reale restano attivi insieme al programma. XDP deve vedere i singoli frame così come arrivano dalla rete, mentre con gli offload di ricezione attivi QEMU consegna alla VM pacchetti già aggregati. Per questo, quando si carica un programma XDP, il driver \texttt{virtio\_net} disattiva automaticamente questi offload; se il dispositivo non permette di disattivarli, rifiuta il programma con l'errore ``Can't set XDP while host is implementing GRO\_HW/CSUM''. È la causa del calo di throughput misurato nel capitolo~\ref{ch:valutazione}, ed è la differenza più evidente rispetto a una SmartNIC reale.
